\section{Experiments}
\label{sec:exp}


\paragraph{Architecture Details.}
Our architecture is identical to Kimi Linear~\citep{zhang2025kimi}, a Mixture-of-Experts (MoE) Transformer following the Moonlight~\citep{liu-2025-moonlight} / DeepSeek-V3~\citep{deepseekaiv3} design, which interleaves Kimi Delta Attention (KDA) and Multi-Head Latent Attention (MLA) layers in a 3:1 ratio, each followed by an MoE feed-forward layer.
The only modification is the addition of AttnRes to the residual connections; all other components (model depth, hidden dimensions, expert routing, and MLP structure) remain unchanged.
AttnRes introduces only one $\operatorname{RMSNorm}$ and one pseudo-query vector $\bm{w}_l \in \mathbb{R}^d$ per layer, amounting to a negligible fraction of the total parameter count. Crucially, all pseudo-query vectors must be initialized to zero. This ensures that the initial attention weights $\brickred{\alpha_{i \to l}}$ are uniform across source layers, which reduces AttnRes to an equal-weight average at the start of training and prevents training volatility, as we validated empirically.


\subsection{Scaling Laws}

\label{sec:scaling}

\begin{table}[t]
    \centering
    \small
    \setlength{\tabcolsep}{6pt}
    \caption{Baseline vs Block AttnRes ($N=8$) vs Full AttnRes vs mHC(-lite) \cite{yang2026mhclite}: Model configurations, Hyperparameters, and Validation Loss.
    }
    \label{tab:baseline_mhc_attnres_compact}
    \resizebox{\textwidth}{!}{%
        \begin{tabular}{cccccccc cccc}
            \toprule
            \multirow{2}{*}{\shortstack{\# Act.                                                                                                                                              \\Params$^\dagger$}} & \multirow{2}{*}{Tokens} & \multirow{2}{*}{$L_b$} & \multirow{2}{*}{$H$} & \multirow{2}{*}{$d_{\text{model}}$} & \multirow{2}{*}{$d_{\text{ff}}$} & \multirow{2}{*}{lr}   & \multirow{2}{*}{batch size$^\ddagger$} &
            \multicolumn{4}{c}{Val.\ Loss}                                                                                                                                                   \\
            \cmidrule(lr){9-12}
                 &                &    &    &              &     &                       &     & Baseline & \shortstack{Block AttnRes} & \shortstack{Full AttnRes} & \shortstack{mHC(-lite)} \\
            \midrule
            194M & \white{0}38.7B & 12 & 12 & \white{0}896 & 400 & $2.99 \times 10^{-3}$ & 192 & 1.931    & 1.909                      & \textbf{1.899}            & 1.906                   \\
            241M & \white{0}45.4B & 13 & 13 & \white{0}960 & 432 & $2.80 \times 10^{-3}$ & 256 & 1.895    & 1.875                      & 1.874                     & \textbf{1.869}          \\
            296M & \white{0}62.1B & 14 & 14 & 1024         & 464 & $2.50 \times 10^{-3}$ & 320 & 1.829    & 1.809                      & \textbf{1.804}            & 1.807                   \\
            436M & \white{0}87.9B & 16 & 16 & 1168         & 528 & $2.20 \times 10^{-3}$ & 384 & 1.766    & 1.746                      & \textbf{1.737}            & 1.747                   \\
            528M & 119.0B         & 17 & 17 & 1264         & 560 & $2.02 \times 10^{-3}$ & 432 & 1.719    & 1.693                      & \textbf{1.692}            & 1.694                   \\
            \bottomrule
            \multicolumn{12}{l}{\rule{0pt}{3ex}$^\dagger$ Denotes the number of activated parameters in our MoE models, excluding embeddings.}                                               \\
            \multicolumn{12}{l}{$^\ddagger$ All models were trained with a context length of 8192.}                                                                                          \\
            \multicolumn{12}{l}{$^\star$ $L_b = L/2$ denotes the number of Transformer blocks.}
        \end{tabular}}
\end{table}

\pgfdeclareplotmark{redstar}{
    \node[star,star point ratio=2.25,minimum size=5pt,
        inner sep=0pt,fill=brickred!80] {};
}
\pgfdeclareplotmark{bluestar}{
    \node[star,star point ratio=2.25,minimum size=5pt,
        inner sep=0pt,fill=midnightblue!80] {};
}
\definecolor{seagreen}{RGB}{235, 165, 85}
\pgfdeclareplotmark{greenstar}{
    \node[star,star point ratio=2.25,minimum size=5pt,
        inner sep=0pt,fill=seagreen] {};
}

\begin{figure}[t]
    \centering
    \begin{tikzpicture}
        \begin{axis}[
                xmode=log,
                ymode=log,
                log ticks with fixed point,
                x filter/.code=\pgfmathparse{#1},
                xlabel={PFLOP/s-days},
                ylabel={Loss},
                xlabel style={font=\small},
                ylabel style={font=\small, at={(0.05,0.5)}},
                xmin=0.5,
                xmax=8,
                ymin=1.68,
                ymax=1.95,
                xtick={0.5, 1, 2, 5},
                xticklabels={0.5, 1, 2, 5},
                ytick={1.7, 1.8, 1.9},
                yticklabels={1.7, 1.8, 1.9},
                width=8cm, height=7.6cm,
                scaled y ticks=false,
                grid=major,
                tick label style={font=\small},
                axis x line=bottom,
                axis y line=left,
                axis line style={-, line width=0.4pt},
                major grid style={line width=0.15pt, gray!30},
                legend style={
                        at={(1,1.03)},
                        anchor=north east,
                        legend cell align=left,
                        font=\scriptsize,
                        draw=none,
                        row sep=-1pt,
                    },
            ]
            \addplot[midnightblue!80, dashed, line width=1.0pt, domain=0.5:8] {
                1.8908 * exp(-0.0565 * ln(x))
            };
            \addlegendentry{Baseline: $1.891 \times C^{-0.057}$}

            \addplot[brickred!80, dashed, line width=1.0pt, domain=0.5:8] {
                1.8645 * exp(-0.0572 * ln(x))
            };
            \addlegendentry{Full AttnRes: $1.865 \times C^{-0.057}$}

            \addplot[seagreen, dashed, line width=1.0pt, domain=0.5:8] {
                1.8699 * exp(-0.0578 * ln(x))
            };
            \addlegendentry{Block AttnRes: $1.870 \times C^{-0.058}$}

            \addplot[color=midnightblue!80, only marks, mark=bluestar] coordinates {
                    (0.6916, 1.931)
                    (1.0405, 1.895)
                    (1.6937, 1.829)
                    (3.3528, 1.766)
                    (5.5974, 1.719)
                };

            \addplot[color=brickred!80, only marks, mark=redstar] coordinates {
                    (0.6916, 1.899)
                    (1.0405, 1.874)
                    (1.6937, 1.804)
                    (3.3528, 1.737)
                    (5.5974, 1.692)
                };

            \addplot[color=seagreen, only marks, mark=greenstar] coordinates {
                    (0.6916, 1.909)
                    (1.0405, 1.875)
                    (1.6937, 1.809)
                    (3.3528, 1.746)
                    (5.5974, 1.693)
                };

            \draw[->, semithick, black] (axis cs:2.39, 1.80) -- (axis cs:1.93, 1.80)
            node[pos=-0.8, above, yshift=2pt, font=\small] {1.25$\times$};

        \end{axis}
    \end{tikzpicture}
    \caption{Scaling law curves for Attention Residuals. Both Full and Block AttnRes consistently outperform the baseline across all scales. Block AttnRes closely tracks Full AttnRes, recovering most of the gain at the largest scale.}
    \label{fig:scaling_law}
\end{figure}


We sweep five model sizes (Table~\ref{tab:baseline_mhc_attnres_compact}) and train three variants per size: a PreNorm baseline, Full AttnRes, and Block AttnRes with ${\approx}\,8$ blocks.
They are trained with an 8192-token context window and a cosine learning rate schedule.
Within each scaling law size group, all variants share identical hyperparameters selected under the baseline to ensure fair comparison; this setup intentionally favors the baseline and thus makes the comparison conservative. Following standard practice, we fit power-law curves of the form $\mathcal{L} = A \times C^{-\alpha}$ \cite{kaplan2020scalinglawsneurallanguage,hoffmann2022trainingcomputeoptimallargelanguage}, where $\mathcal{L}$ is validation loss and $C$ is compute measured in PFLOP/s-days.

\paragraph{Scaling Behavior.}
Fig.~\ref{fig:scaling_law} presents the fitted scaling curves. The Baseline follows $\mathcal{L} = 1.891 \times C^{-0.057}$, while Block AttnRes fits $\mathcal{L} = 1.870 \times C^{-0.058}$, and Full AttnRes fits $\mathcal{L} = 1.865 \times C^{-0.057}$. All three variants exhibit a similar slope, but AttnRes consistently achieves lower loss across the entire compute range. Based on the fitted curves, at 5.6 PFLOP/s-days, Block AttnRes reaches 1.692 versus the Baseline's 1.714, equivalent to a $1.25\times$ compute advantage. The gap between Full and Block AttnRes narrows with scale, shrinking to just 0.001 at the largest size.
We also list mHC(-lite)~\cite{yang2026mhclite} in Table~\ref{tab:baseline_mhc_attnres_compact} for reference. Full AttnRes outperforms mHC, while Block AttnRes matches it at lower memory I/O per layer: $5.5d$ versus $34d$ for mHC with $m{=}4$ streams (Table~\ref{tab:memory_access}).


\subsection{Main Results}
\paragraph{Training recipe.}
The largest models we study are based on the full Kimi Linear 48B configuration: 27 Transformer blocks (54 layers) with 8 out of 256 routed experts plus 1 shared expert, yielding 48B total and 3B activated parameters.
This model applies Block AttnRes with 6 layers per block, producing 9 blocks plus the token embedding for a total of 10 depth-wise sources.

We follow the same data and training recipe as the Kimi Linear 1.4T-token runs~\citep{zhang2025kimi}: all models are pre-trained with a 4096-token context window, the Muon optimizer~\citep{liu-2025-moonlight}, and a WSD (Warmup--Stable--Decay) learning rate schedule~\citep{hu2024minicpm}, with a global batch size of 8M tokens. 
Training of the final model proceeds in two stages: (i) a WSD pre-training phase on 1T tokens, followed by (ii) a mid-training phase on ${\approx}\,$400B high-quality tokens, following the annealing recipe of Moonlight~\citep{liu-2025-moonlight}.

After mid-training, we continue training with progressively longer sequence length of 32K tokens. Since our architecture uses hybrid KDA/MLA attention~\citep{zhang2025kimi}, where MLA operates without positional encodings (NoPE)~\cite{yang2025ropenopeagainnew}, context extension requires no modifications such as YaRN~\citep{peng2023yarn} or attention temperature rescaling.

\label{sec:main}

\begin{figure*}[h]
    \centering
    \small
    \pgfplotsset{
        every axis/.style={
                width=0.85\linewidth,
                height=120pt,
                scale only axis,
                grid=major,
                grid style={black!15, thin},
                axis x line*=bottom,
                x axis line style={line width=0.8pt},
                axis y line*=left,
                y axis line style={draw=none},
                every y tick/.style={draw=none},
                xlabel style={font=\scriptsize},
                ylabel=,
                title style={font=\scriptsize, at={(0.5,1)}, anchor=south},
                xticklabel style={font=\scriptsize},
                yticklabel style={font=\scriptsize},
                legend style={font=\scriptsize, fill=none, legend cell align=left},
                scaled x ticks=false,
                enlargelimits=false,
                clip=true,
            },
        baseline/.style={midnightblue!80, line width=0.8pt, mark=hexagon*, mark size=0.8pt, mark options={fill=midnightblue!80}},
        attnres/.style={brickred!80, line width=0.8pt, mark=hexagon*, mark size=0.8pt, mark options={fill=brickred!80}},
    }
    \begin{subfigure}[t]{0.33\textwidth}
        \begin{tikzpicture}[baseline=(loss.north)]
            \begin{axis}[
                    name=loss,
                    title={(a) Validation Loss},
                    xlabel={Step},
                    xmin=10000, xmax=120000,
                    ymin=1.14, ymax=1.50,
                    xtick={20000,40000,60000,80000,100000},
                    xticklabels={20k,40k,60k,80k,100k},
                    legend style={at={(1.0,1.0)}, anchor=north east},
                ]
                \addplot[baseline, mark=none] coordinates {
                        (10000,1.4827) (11000,1.4778) (12000,1.4724) (13000,1.4666) (14000,1.4604) (15000,1.4542)
                        (16000,1.4479) (17000,1.4419) (18000,1.4361) (19000,1.4307) (20000,1.4257) (21000,1.4211)
                        (22000,1.4169) (23000,1.4132) (24000,1.4097) (25000,1.4066) (26000,1.4037) (27000,1.401)
                        (28000,1.3985) (29000,1.396) (30000,1.3936) (31000,1.3913) (32000,1.389) (33000,1.3867)
                        (34000,1.3845) (35000,1.3823) (36000,1.3802) (37000,1.3781) (38000,1.3762) (39000,1.3744)
                        (40000,1.3727) (41000,1.3711) (42000,1.3696) (43000,1.3682) (44000,1.3668) (45000,1.3655)
                        (46000,1.3642) (47000,1.363) (48000,1.3617) (49000,1.3605) (50000,1.3594) (51000,1.3583)
                        (52000,1.3573) (53000,1.3564) (54000,1.3556) (55000,1.3548) (56000,1.3541) (57000,1.3534)
                        (58000,1.3527) (59000,1.352) (60000,1.3512) (61000,1.3502) (62000,1.3493) (63000,1.3482)
                        (64000,1.3471) (65000,1.3461) (66000,1.345) (67000,1.3441) (68000,1.3432) (69000,1.3425)
                        (70000,1.3419) (71000,1.3413) (72000,1.3408) (73000,1.3404) (74000,1.3399) (75000,1.3395)
                        (76000,1.339) (77000,1.3384) (78000,1.3379) (79000,1.3373) (80000,1.3367) (81000,1.336)
                        (82000,1.3353) (83000,1.3345) (84000,1.3335) (85000,1.3325) (86000,1.3312) (87000,1.3298)
                        (88000,1.3281) (89000,1.3261) (90000,1.3239) (91000,1.3214) (92000,1.3186) (93000,1.3154)
                        (94000,1.3119) (95000,1.308) (96000,1.3038) (97000,1.2992) (98000,1.2941) (99000,1.2888)
                        (100000,1.2831) (101000,1.2771) (102000,1.2708) (103000,1.2644) (104000,1.2578) (105000,1.2512)
                        (106000,1.2444) (107000,1.2377) (108000,1.231) (109000,1.2244) (110000,1.218) (111000,1.2117)
                        (112000,1.2058) (113000,1.2002) (114000,1.1949) (115000,1.1902) (116000,1.1858) (117000,1.182)
                        (118000,1.1787) (119000,1.1758) (119209,1.1734)
                    };
                \addplot[attnres, mark=none] coordinates {
                        (10000,1.4736) (11000,1.4688) (12000,1.4635) (13000,1.4576) (14000,1.4515) (15000,1.4452)
                        (16000,1.4389) (17000,1.4327) (18000,1.4267) (19000,1.4212) (20000,1.416) (21000,1.4113)
                        (22000,1.407) (23000,1.4031) (24000,1.3996) (25000,1.3963) (26000,1.3933) (27000,1.3905)
                        (28000,1.3879) (29000,1.3854) (30000,1.3829) (31000,1.3805) (32000,1.3782) (33000,1.3759)
                        (34000,1.3737) (35000,1.3715) (36000,1.3694) (37000,1.3674) (38000,1.3655) (39000,1.3636)
                        (40000,1.3618) (41000,1.3602) (42000,1.3585) (43000,1.357) (44000,1.3554) (45000,1.3539)
                        (46000,1.3525) (47000,1.3511) (48000,1.3497) (49000,1.3483) (50000,1.347) (51000,1.3458)
                        (52000,1.3447) (53000,1.3436) (54000,1.3427) (55000,1.3419) (56000,1.341) (57000,1.3403)
                        (58000,1.3394) (59000,1.3386) (60000,1.3376) (61000,1.3366) (62000,1.3355) (63000,1.3343)
                        (64000,1.3331) (65000,1.332) (66000,1.3308) (67000,1.3298) (68000,1.3289) (69000,1.328)
                        (70000,1.3273) (71000,1.3267) (72000,1.3261) (73000,1.3255) (74000,1.325) (75000,1.3244)
                        (76000,1.3238) (77000,1.3232) (78000,1.3226) (79000,1.322) (80000,1.3213) (81000,1.3207)
                        (82000,1.32) (83000,1.3192) (84000,1.3183) (85000,1.3172) (86000,1.3159) (87000,1.3145)
                        (88000,1.3127) (89000,1.3107) (90000,1.3084) (91000,1.3058) (92000,1.3029) (93000,1.2996)
                        (94000,1.296) (95000,1.292) (96000,1.2877) (97000,1.2829) (98000,1.2778) (99000,1.2723)
                        (100000,1.2665) (101000,1.2604) (102000,1.254) (103000,1.2475) (104000,1.2408) (105000,1.234)
                        (106000,1.2272) (107000,1.2203) (108000,1.2135) (109000,1.2068) (110000,1.2002) (111000,1.1939)
                        (112000,1.1878) (113000,1.1821) (114000,1.1768) (115000,1.1719) (116000,1.1675) (117000,1.1637)
                        (118000,1.1604) (119000,1.1575) (119209,1.1551)
                    };
                \legend{Baseline, Block AttnRes}
            \end{axis}
            \draw[black, line width=0.8pt] (loss.north west) -- (loss.north east);
        \end{tikzpicture}
        \phantomcaption\label{fig:comparison-loss}
    \end{subfigure}%
    \hfill%
    \begin{subfigure}[t]{0.33\textwidth}
        \begin{tikzpicture}[baseline=(rms.north)]
            \begin{axis}[
                    name=rms,
                    title={(b) Output Magnitude},
                    xlabel={Transformer Block Index},
                    xmin=-0.5, xmax=26.5,
                    ymin=0, ymax=15,
                    ytick={0,5,10,15},
                ]
                \addplot[baseline, mark=none] coordinates {
                        (0,0.04) (1,0.06) (2,0.10) (3,0.14) (4,0.20) (5,0.59) (6,0.63) (7,0.68) (8,0.71)
                        (9,0.76) (10,0.81) (11,0.89) (12,0.98) (13,1.07) (14,1.23) (15,1.43) (16,1.66) (17,1.91)
                        (18,2.29) (19,2.60) (20,3.12) (21,3.51) (22,5.56) (23,6.20) (24,7.10) (25,10.47) (26,12.15)
                    };
                \addplot[attnres, mark=none] coordinates {
                        (0,0.07) (1,0.34) (2,0.75) (3,0.35) (4,0.67) (5,1.89) (6,0.38) (7,0.69) (8,1.30)
                        (9,0.30) (10,0.66) (11,1.48) (12,0.21) (13,0.42) (14,0.71) (15,0.28) (16,0.59) (17,1.48)
                        (18,0.46) (19,0.75) (20,1.41) (21,0.66) (22,1.23) (23,1.50) (24,0.64) (25,1.16) (26,1.91)
                    };
            \end{axis}
            \draw[black, line width=0.8pt] (rms.north west) -- (rms.north east);
        \end{tikzpicture}
        \phantomcaption\label{fig:comparison-rms}
    \end{subfigure}%
    \hfill%
    \begin{subfigure}[t]{0.33\textwidth}
        \begin{tikzpicture}[baseline=(grad.north)]
            \begin{axis}[
                    name=grad,
                    title={(c) Gradient Magnitude ($\times 10^{-5}$)},
                    xlabel={Transformer Block Index},
                    xmin=-0.5, xmax=26.5,
                    ymin=0, ymax=3,
                    ytick={0,1,2,3},
                ]
                \addplot[baseline, mark=none] coordinates {
                        (0,2.4663) (1,1.8375) (2,1.3963) (3,1.1747) (4,0.9401) (5,0.7749) (6,0.6823)
                        (7,0.5613) (8,0.4734) (9,0.4208) (10,0.3728) (11,0.3159) (12,0.2787) (13,0.2485)
                        (14,0.2230) (15,0.1892) (16,0.1636) (17,0.1464) (18,0.1340) (19,0.1146) (20,0.1011)
                        (21,0.0904) (22,0.0815) (23,0.0737) (24,0.0684) (25,0.0698) (26,0.0844)
                    };
                \addplot[attnres, mark=none] coordinates {
                        (0,0.7177) (1,0.3159) (2,0.2428) (3,0.2136) (4,0.1559) (5,0.1346) (6,0.1797)
                        (7,0.1404) (8,0.1116) (9,0.2345) (10,0.1444) (11,0.0869) (12,0.3832) (13,0.2231)
                        (14,0.1624) (15,0.2936) (16,0.1630) (17,0.0959) (18,0.2454) (19,0.1708) (20,0.1221)
                        (21,0.1970) (22,0.1694) (23,0.1553) (24,0.2260) (25,0.2237) (26,0.2859)
                    };
            \end{axis}
            \draw[black, line width=0.8pt] (grad.north west) -- (grad.north east);
        \end{tikzpicture}
        \phantomcaption\label{fig:comparison-grad}
    \end{subfigure}

    \caption{Training dynamics of Baseline and Block AttnRes. \textbf{(\subref{fig:comparison-loss})} Validation loss during training. \textbf{(\subref{fig:comparison-rms})} Each transformer block's output magnitude at the end of training. \textbf{(\subref{fig:comparison-grad})} Each transformer block's gradient magnitude.
    }
    \label{fig:comparison}
\end{figure*}

\paragraph{Training dynamics.}
We compare the training dynamics of our final Baseline and Block AttnRes models over 1T tokens in Fig.~\ref{fig:comparison}.
\begin{itemize}[leftmargin=*]
    \item \textbf{Validation loss:} AttnRes achieves consistently lower validation loss throughout training, with the gap widening during the decay phase and resulting in a notably lower final loss.
    \item \textbf{Output magnitude:} The Baseline suffers from the PreNorm dilution problem~\citep{xiong2020layer,li2026siamesenorm}: as hidden-state magnitudes grow monotonically with depth, deeper layers are compelled to learn increasingly large outputs from fixed-scale normalized inputs to remain influential. Block AttnRes confines this growth within each block, as selective aggregation at block boundaries resets the accumulation, yielding a bounded periodic pattern.
    \item \textbf{Gradient magnitude:} With all residual weights fixed to 1, the Baseline provides no means of regulating gradient flow across depth, leading to disproportionately large gradients in the earliest layers. The learnable $\operatorname{softmax}$ weights in Block AttnRes (Fig.~\ref{fig:attn-res-weights}) introduce competition among sources for probability mass, resulting in a substantially more uniform gradient distribution.
\end{itemize}

\begin{table}[!ht]
    \centering
    \small
    \renewcommand{\arraystretch}{1.1}
    \caption{Performance comparison of AttnRes with the baseline, both after the same pre-training recipe. Best per-row results are \textbf{bolded}.}
    \setlength{\tabcolsep}{12pt}
    \label{tab:benchmarkresult}
    \begin{tabular}{@{}r l c c}
        \toprule
         &               & Baseline      & AttnRes       \\
        \midrule
        \multirow{7}{*}{\textit{General}}
         & MMLU          & 73.5          & \textbf{74.6} \\
         & MMLU-Pro      & \textbf{52.2} & \textbf{52.2} \\
         & GPQA-Diamond  & 36.9          & \textbf{44.4} \\
         & BBH           & 76.3          & \textbf{78.0} \\
         & ARC-Challenge & 64.6          & \textbf{65.7} \\
         & HellaSwag     & 83.2          & \textbf{83.4} \\
         & TriviaQA      & 69.9          & \textbf{71.8} \\
        \midrule
        \multirow{6}{*}{\textit{Math \& Code}}
         & GSM8K         & 81.7          & \textbf{82.4} \\
         & MGSM          & 64.9          & \textbf{66.1} \\
         & Math          & 53.5          & \textbf{57.1} \\
         & CMath         & 84.7          & \textbf{85.1} \\
         & HumanEval     & 59.1          & \textbf{62.2} \\
         & MBPP          & 72.0          & \textbf{73.9} \\
        \midrule
        \multirow{2}{*}{\textit{Chinese}}
         & CMMLU         & 82.0          & \textbf{82.9} \\
         & C-Eval        & 79.6          & \textbf{82.5} \\
        \bottomrule
    \end{tabular}
    \label{tab:pretrain}
\end{table}

\paragraph{Downstream performance.}
Following the evaluation protocol of Kimi Linear~\cite{zhang2025kimi}, we assess both models across three areas (Table~\ref{tab:pretrain}):
\begin{itemize}[leftmargin=*]
    \item \textbf{Language understanding and reasoning}: MMLU~\cite{hendrycks-2021-mmlu}, MMLU-Pro Hard~\cite{wang2024mmlu}, GPQA-Diamond~\cite{rein2024gpqa}, BBH~\cite{suzgun2022challenging}, ARC-Challenge~\cite{clark-2018-arc}, HellaSwag~\cite{zellers-2019-hellaswag}, and TriviaQA~\cite{joshi2017triviaqa}.
    \item \textbf{Reasoning (Code and Math)}: GSM8K~\cite{gsm8k}, MGSM~\cite{mgsm}, Math~\cite{minervamath}, CMath~\cite{cmathbmk}, HumanEval~\cite{chen2021evaluatinglargelanguagemodels}, and MBPP~\cite{austin2021programsynthesislargelanguage}.
    \item \textbf{Chinese language understanding}: CMMLU~\cite{li-etal-2024-cmmlu} and C-Eval~\cite{huang2023c}.
\end{itemize}
As shown in Table~\ref{tab:pretrain}, Block AttnRes matches or outperforms the baseline on all benchmarks. The improvements are particularly pronounced on multi-step reasoning tasks such as GPQA-Diamond (+7.5) and Minerva Math (+3.6), as well as code generation such as HumanEval (+3.1), while knowledge-oriented benchmarks such as MMLU (+1.1) and TriviaQA (+1.9) also show solid gains. This pattern is consistent with the hypothesis that improved depth-wise information flow benefits compositional tasks, where later layers can selectively retrieve and build upon earlier representations.

\subsection{Ablation Study}
\pgfdeclareplotmark{hexagon*}{
    \pgfpathmoveto{\pgfqpointpolar{0}{1.1\pgfplotmarksize}}
    \pgfpathlineto{\pgfqpointpolar{60}{1.1\pgfplotmarksize}}
    \pgfpathlineto{\pgfqpointpolar{120}{1.1\pgfplotmarksize}}
    \pgfpathlineto{\pgfqpointpolar{180}{1.1\pgfplotmarksize}}
    \pgfpathlineto{\pgfqpointpolar{240}{1.1\pgfplotmarksize}}
    \pgfpathlineto{\pgfqpointpolar{300}{1.1\pgfplotmarksize}}
    \pgfpathclose
    \pgfusepathqfillstroke
}
\begin{figure}[t!]
    \centering
    \adjustbox{valign=t}{\begin{minipage}{0.45\textwidth}
            \centering
            \small
            \setlength{\tabcolsep}{10pt}
            \renewcommand{\arraystretch}{1.18}
            \captionof{table}{Ablation on key components of AttnRes (16-layer model).}
            \label{tab:ablation}
            \begin{tabular}{llc}
                \toprule
                                                                                  & Variant                                   & Loss  \\
                \midrule
                \multicolumn{2}{l}{Baseline (PreNorm)}                            & 1.766                                             \\
                \midrule
                \multicolumn{2}{l}{DenseFormer \cite{pagliardini2024denseformer}} & 1.767                                             \\
                \multicolumn{2}{l}{mHC \cite{xie2026mhc}}                         & 1.747                                             \\
                \midrule
                AttnRes                                                           & Full                                      & 1.737 \\
                                                                                  & \quad \emph{w/ input-dependent query}     & 1.731 \\
                                                                                  & \quad \emph{w/ input-independent mixing}  & 1.749 \\
                                                                                  & \quad \emph{w/ $\operatorname{sigmoid}$}  & 1.741 \\
                                                                                  & \quad \emph{w/o $\operatorname{RMSNorm}$} & 1.743 \\
                                                                                  & SWA ($W = 1 + 8$)                         & 1.764 \\
                                                                                  & Block ($S=4$)                             & 1.746 \\
                                                                                  & \quad \emph{w/ multihead ($H=16$)}        & 1.752 \\
                                                                                  & \quad \emph{w/o $\operatorname{RMSNorm}$} & 1.750 \\
                \bottomrule
            \end{tabular}
        \end{minipage}}
    \hfill
    \adjustbox{valign=t}{\begin{minipage}{0.5\textwidth}
            \centering
            \small
            \begin{tikzpicture}
                \begin{axis}[
                        xlabel={Block size ($S$)},
                        ylabel={Validation loss},
                        xlabel style={font=\small},
                        ylabel style={font=\small},
                        separate axis lines,
                        every outer x axis line/.append style={line width=0.8pt},
                        every outer y axis line/.append style={draw=none},
                        every y tick/.style={draw=none},
                        every x tick/.style={draw=none},
                        xticklabel style={font=\small},
                        yticklabel style={font=\small, /pgf/number format/fixed, /pgf/number format/precision=3, /pgf/number format/zerofill},
                        xtick={1, 2, 3, 4, 5},
                        xticklabels={32, 16, 8, 4, 2},
                        xmin=0.5,
                        xmax=5.5,
                        ymin=1.734,
                        ymax=1.770,
                        ytick={1.735, 1.740, 1.745, 1.750, 1.755, 1.760, 1.765, 1.770},
                        scaled y ticks=false,
                        width=\textwidth,
                        height=7cm,
                        grid=major,
                        grid style={gray!20, thin},
                        enlargelimits=false,
                        legend style={
                                at={(1.,1.)},
                                anchor=north east,
                                legend cell align=left,
                                font=\scriptsize,
                                draw=gray!50,
                                fill=white,
                                fill opacity=0.9,
                                text opacity=1,
                            },
                    ]

                    \addplot[gray!60, dashed, line width=0.8pt, forget plot] coordinates {(0.5, 1.766) (5.5, 1.766)};
                    \addlegendimage{gray!60, dashed, line width=0.8pt}
                    \addlegendentry{Baseline}

                    \addplot[brickred!80, dashed, line width=0.8pt, forget plot] coordinates {(0.5, 1.737) (5.5, 1.737)};
                    \addlegendimage{brickred!80, dashed, line width=0.8pt}
                    \addlegendentry{Full AttnRes}

                    \addplot[brickred!80, line width=0.8pt, mark=hexagon*, mark size=1.2pt, mark options={fill=brickred!80}] coordinates {
                            (1, 1.757)
                            (2, 1.753)
                            (3, 1.748)
                            (4, 1.746)
                            (5, 1.746)
                        };
                    \addlegendentry{Block AttnRes}

                    \node[font=\tiny, above=3pt] at (axis cs:1, 1.757) {1.757};
                    \node[font=\tiny, above=3pt] at (axis cs:2, 1.753) {1.753};
                    \node[font=\tiny, above=3pt] at (axis cs:3, 1.748) {1.748};
                    \node[font=\tiny, above=3pt] at (axis cs:4, 1.746) {1.746};
                    \node[font=\tiny, above=3pt] at (axis cs:5, 1.746) {1.746};

                    \node[font=\tiny, black, above=2pt] at (axis cs:3, 1.766) {Baseline (1.766)};

                    \node[font=\tiny, black, below=2pt,text=brickred!80] at (axis cs:3, 1.737) {Full AttnRes i.e. S=1 (1.737)};

                \end{axis}
            \end{tikzpicture}
            \captionof{figure}{Effect of block size on validation loss (16-layer model).}
            \label{fig:blocksize}
        \end{minipage}}
\end{figure}

We conduct ablation studies on the 16-head model from Table~\ref{tab:baseline_mhc_attnres_compact} to validate key design choices in AttnRes (Table~\ref{tab:ablation}). All models share identical hyperparameters and compute budget.

\paragraph{Comparison with prior methods.}
We compare AttnRes against the PreNorm baseline (loss 1.766) and two representative methods that generalize residual connections. DenseFormer~\cite{pagliardini2024denseformer} grants each layer access to all previous outputs but combines them with fixed, input-independent scalar coefficients; it shows no gain over the baseline (1.767), highlighting the importance of input-dependent weighting. mHC~\cite{xie2026mhc} introduces input dependence through $m$ parallel streams with learned mixing matrices, improving to 1.747. AttnRes takes this further with explicit content-dependent selection via $\operatorname{softmax}$ attention: Full AttnRes achieves 1.737 and Block AttnRes 1.746, outperforming both methods with only a single query vector per layer.

\paragraph{Cross-layer access.}
We compare three granularities of cross-layer access. Full AttnRes follows directly from the time--depth duality (\S~\ref{sec:attnres}), applying attention over all previous layers, and achieves the lowest loss (1.737). A simple way to reduce its memory cost is sliding-window aggregation (SWA), which retains only the most recent $W{=}8$ layer outputs plus the token embedding; it improves over baseline (1.764) but falls well short of both Full and Block AttnRes, suggesting that selectively accessing distant layers matters more than attending to many nearby ones.

Block AttnRes offers a better trade-off: with block size $S{=}4$ it reaches 1.746 while keeping memory overhead constant per layer. Fig.~\ref{fig:blocksize} sweeps $S$ across the full spectrum from $S{=}1$ (i.e.\ Full AttnRes) to increasingly coarse groupings. Loss degrades gracefully as $S$ grows, with $S{=}2, 4, 8$ all landing near 1.746 while larger blocks ($S{=}16, 32$) move toward baseline. In practice, we fix the number of blocks to ${\approx}\,8$ for infrastructure efficiency (\S~\ref{sec:infra}). As future hardware alleviates memory capacity constraints, adopting finer-grained block sizes or Full AttnRes represents a natural pathway to further improve performance.

\paragraph{Component design.}
We further ablate individual components of the attention mechanism:
\begin{itemize}[leftmargin=*]
    \item \textbf{Input-dependent query.} A natural extension is to make the query input-dependent by projecting it from the current hidden state. This further lowers loss to 1.731, but introduces a $d \times d$ projection per layer and requires sequential memory access during decoding, so we default to the learned query.
    \item \textbf{Input-independent mixing.} We removed the query and key and replaced them with learnable, input-independent scalars to weigh previous layers, which hurts performance (1.749 vs.\ 1.737).
    \item \textbf{$\operatorname{softmax}$ vs.\ $\operatorname{sigmoid}$.} Replacing $\operatorname{softmax}$ with $\operatorname{sigmoid}$ degrades performance (1.741). We attribute this to $\operatorname{softmax}$'s competitive normalization, which forces sharper selection among sources.
    \item \textbf{Multihead attention.} We test per-head depth aggregation ($H{=}16$) on Block AttnRes, allowing different channel groups to attend to different source layers. This hurts performance (1.752 vs.\ 1.746), indicating that the optimal depth-wise mixture is largely uniform across channels: when a layer's output is relevant, it is relevant as a whole.
    \item \textbf{$\operatorname{RMSNorm}$ on keys.} Removing $\operatorname{RMSNorm}$ degrades both Full AttnRes (1.743) and Block AttnRes (1.750). For Full AttnRes, it prevents individual layers with naturally larger outputs from dominating the $\operatorname{softmax}$. This becomes even more critical for Block AttnRes, as block-level representations accumulate over more layers and can develop large magnitude differences; $\operatorname{RMSNorm}$ prevents these from biasing the attention weights.
\end{itemize}

\subsection{Analysis}



\subsubsection{Optimal Architecture}
\label{sec:optimal_arch}

\begin{figure}[t]
  \centering
  \begin{subfigure}[b]{0.48\textwidth}
    \centering
    \begin{tikzpicture}
      \begin{axis}[
          width=\textwidth,
          height=0.85\textwidth,
          xlabel={$d_{\mathrm{model}}/L_b$},
          ylabel={$H/L_b$},
          xtick={15,30,45,60,75},
          ytick={0.3,0.4,0.5,0.6,0.7},
          xticklabel style={font=\small},
          yticklabel style={font=\small},
          xlabel style={font=\small},
          ylabel style={font=\small, at={(0.04,0.5)}},
          colormap={softcoolwarm}{
              rgb255(0)=(60,110,190)
              rgb255(1)=(100,150,205)
              rgb255(2)=(150,185,215)
              rgb255(3)=(205,210,215)
              rgb255(4)=(225,195,175)
              rgb255(5)=(220,150,110)
              rgb255(6)=(210,100,75)
              rgb255(7)=(190,60,50)
            },
          title style={font=\small},
          point meta min=1.80,
          point meta max=2.04,
          enlargelimits=false,
        ]
        \addplot[matrix plot*, mesh/cols=5, point meta=explicit] coordinates {
            (15,0.7) [2.017] (30,0.7) [1.909] (45,0.7) [1.875] (60,0.7) [1.851] (75,0.7) [1.858]
            (15,0.6) [1.990] (30,0.6) [1.902] (45,0.6) [1.862] (60,0.6) [1.852] (75,0.6) [1.862]
            (15,0.5) [1.973] (30,0.5) [1.883] (45,0.5) [1.859] (60,0.5) [1.849] (75,0.5) [1.854]
            (15,0.4) [1.952] (30,0.4) [1.868] (45,0.4) [1.850] (60,0.4) [1.849] (75,0.4) [1.857]
            (15,0.3) [1.926] (30,0.3) [1.857] (45,0.3) [1.851] (60,0.3) [1.847] (75,0.3) [1.858]
          };

        \node[font=\scriptsize, text=white] at (axis cs:15,0.7) {2.017};
        \node[font=\scriptsize, text=black] at (axis cs:30,0.7) {1.909};
        \node[font=\scriptsize, text=black] at (axis cs:45,0.7) {1.875};
        \node[font=\scriptsize, text=black] at (axis cs:60,0.7) {1.851};
        \node[font=\scriptsize, text=black] at (axis cs:75,0.7) {1.858};
        \node[font=\scriptsize, text=white] at (axis cs:15,0.6) {1.990};
        \node[font=\scriptsize, text=black] at (axis cs:30,0.6) {1.902};
        \node[font=\scriptsize, text=black] at (axis cs:45,0.6) {1.862};
        \node[font=\scriptsize, text=black] at (axis cs:60,0.6) {1.852};
        \node[font=\scriptsize, text=black] at (axis cs:75,0.6) {1.862};
        \node[font=\scriptsize, text=white] at (axis cs:15,0.5) {1.973};
        \node[font=\scriptsize, text=black] at (axis cs:30,0.5) {1.883};
        \node[font=\scriptsize, text=black] at (axis cs:45,0.5) {1.859};
        \node[font=\scriptsize, text=black] at (axis cs:60,0.5) {1.849};
        \node[font=\scriptsize, text=black] at (axis cs:75,0.5) {1.854};
        \node[font=\scriptsize, text=white] at (axis cs:15,0.4) {1.952};
        \node[font=\scriptsize, text=black] at (axis cs:30,0.4) {1.868};
        \node[font=\scriptsize, text=black] at (axis cs:45,0.4) {1.850};
        \node[font=\scriptsize, text=black] at (axis cs:60,0.4) {1.849};
        \node[font=\scriptsize, text=black] at (axis cs:75,0.4) {1.857};
        \node[font=\scriptsize, text=black] at (axis cs:15,0.3) {1.926};
        \node[font=\scriptsize, text=black] at (axis cs:30,0.3) {1.857};
        \node[font=\scriptsize, text=black] at (axis cs:45,0.3) {1.851};
        \node[font=\scriptsize, text=black] at (axis cs:75,0.3) {1.858};

        \node[star, star point ratio=2.25, minimum size=10pt,
          inner sep=0pt, draw=brickred, solid, fill=brickred]
        at (axis cs:60,0.3) {};
        \node[font=\scriptsize, text=black] at (axis cs:60,0.3) {1.847};
      \end{axis}
    \end{tikzpicture}
    \caption{Baseline}
    \label{fig:arch_sweep_baseline}
  \end{subfigure}
  \hfill
  \begin{subfigure}[b]{0.48\textwidth}
    \centering
    \begin{tikzpicture}
      \begin{axis}[
          width=\textwidth,
          height=0.85\textwidth,
          xlabel={$d_{\mathrm{model}}/L_b$},
          xtick={15,30,45,60,75},
          xticklabel style={font=\small},
          yticklabel style={font=\small},
          yticklabels={,,},
          xlabel style={font=\small},
          ylabel style={font=\small, at={(0.04,0.5)}},
          colormap={softcoolwarm}{
              rgb255(0)=(60,110,190)
              rgb255(1)=(100,150,205)
              rgb255(2)=(150,185,215)
              rgb255(3)=(205,210,215)
              rgb255(4)=(225,195,175)
              rgb255(5)=(220,150,110)
              rgb255(6)=(210,100,75)
              rgb255(7)=(190,60,50)
            },
          colorbar,
          colorbar style={
              width=6pt,
              xshift=4pt,
              ytick={1.84, 1.88, 1.92, 1.96, 2.00},
              yticklabel style={font=\scriptsize},
            },
          title style={font=\small},
          point meta min=1.80,
          point meta max=2.04,
          enlargelimits=false,
        ]
        \addplot[matrix plot*, mesh/cols=5, point meta=explicit] coordinates {
            (15,0.7) [1.954] (30,0.7) [1.890] (45,0.7) [1.843] (60,0.7) [1.828] (75,0.7) [1.824]
            (15,0.6) [1.931] (30,0.6) [1.863] (45,0.6) [1.830] (60,0.6) [1.817] (75,0.6) [1.818]
            (15,0.5) [1.917] (30,0.5) [1.841] (45,0.5) [1.819] (60,0.5) [1.812] (75,0.5) [1.817]
            (15,0.4) [1.893] (30,0.4) [1.823] (45,0.4) [1.815] (60,0.4) [1.813] (75,0.4) [1.813]
            (15,0.3) [1.877] (30,0.3) [1.816] (45,0.3) [1.802] (60,0.3) [1.806] (75,0.3) [1.820]
          };

        \node[font=\scriptsize, text=white] at (axis cs:15,0.7) {1.954};
        \node[font=\scriptsize, text=black] at (axis cs:30,0.7) {1.890};
        \node[font=\scriptsize, text=black] at (axis cs:45,0.7) {1.843};
        \node[font=\scriptsize, text=black] at (axis cs:60,0.7) {1.828};
        \node[font=\scriptsize, text=black] at (axis cs:75,0.7) {1.824};
        \node[font=\scriptsize, text=black] at (axis cs:15,0.6) {1.931};
        \node[font=\scriptsize, text=black] at (axis cs:30,0.6) {1.863};
        \node[font=\scriptsize, text=black] at (axis cs:45,0.6) {1.830};
        \node[font=\scriptsize, text=black] at (axis cs:60,0.6) {1.817};
        \node[font=\scriptsize, text=black] at (axis cs:75,0.6) {1.818};
        \node[font=\scriptsize, text=black] at (axis cs:15,0.5) {1.917};
        \node[font=\scriptsize, text=black] at (axis cs:30,0.5) {1.841};
        \node[font=\scriptsize, text=black] at (axis cs:45,0.5) {1.819};
        \node[font=\scriptsize, text=black] at (axis cs:60,0.5) {1.812};
        \node[font=\scriptsize, text=black] at (axis cs:75,0.5) {1.817};
        \node[font=\scriptsize, text=black] at (axis cs:15,0.4) {1.893};
        \node[font=\scriptsize, text=black] at (axis cs:30,0.4) {1.823};
        \node[font=\scriptsize, text=black] at (axis cs:45,0.4) {1.815};
        \node[font=\scriptsize, text=black] at (axis cs:60,0.4) {1.813};
        \node[font=\scriptsize, text=black] at (axis cs:75,0.4) {1.813};
        \node[font=\scriptsize, text=black] at (axis cs:15,0.3) {1.877};
        \node[font=\scriptsize, text=black] at (axis cs:30,0.3) {1.816};
        \node[font=\scriptsize, text=black] at (axis cs:75,0.3) {1.820};
        \node[font=\scriptsize, text=black] at (axis cs:60,0.3) {1.806};

        \node[star, star point ratio=2.25, minimum size=10pt,
          inner sep=0pt, draw=brickred, solid, fill=brickred]
        at (axis cs:45,0.3) {};
        \node[font=\scriptsize, text=black] at (axis cs:45,0.3) {1.802};

      \end{axis}
    \end{tikzpicture}
    \caption{Attention Residuals}
    \label{fig:arch_sweep_attnres}
  \end{subfigure}
  \caption{
    Architecture sweep under fixed compute ($\approx 6.5\times 10^{19}$ FLOPs, $\approx 2.3\times 10^{8}$ active parameters).
    Each cell reports validation loss for a $(d_{\mathrm{model}}/L_b,\; H/L_b)$ configuration, where $L_b = L/2$ is the number of Transformer blocks; the star marks the optimum.
  }
  \label{fig:arch_sweep}
\end{figure}



To understand how AttnRes reshapes optimal architectural scaling, we perform a controlled capacity reallocation study under a fixed compute and parameter budget. Our central question is whether AttnRes alters the preferred depth–width–attention trade-off, and in particular, given its potential strength on the depth dimension, whether it favors deeper models compared to conventional Transformer design heuristics. To isolate structural factors directly coupled to depth, we fix the per-expert MLP expansion ratio based on internal empirical observations ($d_{\mathrm{ff}}/d_{\mathrm{model}} \approx 0.45$). We further fix total training compute (FLOPs $\approx 6.5\times10^{19}$) and active parameters ($\approx 2.3\times10^8$), ensuring that any performance variation arises purely from architectural reallocation rather than overall capacity differences. Under this constrained budget, we enumerate 25 configurations on a $5\times5$ grid over
$d_{\mathrm{model}}/L_b \in \{15, 30, 45, 60, 75\}$ and
$H/L_b \in \{0.3, 0.4, 0.5, 0.6, 0.7\}$,
where $L_b = L/2$ is the number of Transformer blocks and $H$ the number of attention heads.
The results are shown in Fig.~\ref{fig:arch_sweep}.

Both heatmaps exhibit a shared pattern: loss decreases with growing $d_{\mathrm{model}}/L_b$ and shrinking $H/L_b$, and both methods reach their optima at $H/L_b \approx 0.3$.
Despite this shared trend, AttnRes achieves a lower loss than the baseline in each of the 25 configurations, by $0.019$--$0.063$.
The most apparent difference lies in the location of the optimum: the baseline achieves its lowest loss at $d_{\mathrm{model}}/L_b \approx 60$ ($1.847$), whereas AttnRes shifts it to $d_{\mathrm{model}}/L_b \approx 45$ ($1.802$).
Under a fixed parameter budget, a lower $d_{\mathrm{model}}/L_b$ corresponds to a deeper, narrower network, suggesting that AttnRes can exploit additional depth more effectively.
We note that this preference for depth does not directly translate to a deployment recommendation, as deeper models generally incur higher inference latency due to their sequential computation~\citep{pope-2022-efficiently}.
Rather, this sweep serves as a diagnostic that reveals where AttnRes benefits most, and this depth preference can be factored into the architecture selection alongside inference cost.


\subsubsection{Analyzing Learned AttnRes Patterns}

\begin{figure*}[t]
  \centering
  \adjustbox{width=\textwidth}{%
    \begin{tikzpicture}
      \pgfdeclarepatternformonly{sparse north east lines}
        {\pgfqpoint{-0.5pt}{-0.5pt}}{\pgfqpoint{9.5pt}{9.5pt}}{\pgfqpoint{9pt}{9pt}}
        {\pgfsetlinewidth{0.4pt}\pgfpathmoveto{\pgfqpoint{0pt}{0pt}}\pgfpathlineto{\pgfqpoint{9pt}{9pt}}\pgfusepath{stroke}}
      \pgfplotsset{
        heatbase/.style={
            every axis title/.style={at={(0.5,1.05)}, anchor=south, font=\small},
            xlabel style={font=\small},
            ylabel style={font=\scriptsize, at={(0.05,0.5)}},
            xticklabel style={font=\small},
            yticklabel style={font=\small},
            colormap={whiteblue}{
                rgb255(0cm)=(255,255,255)
                rgb255(1cm)=(220,235,248)
                rgb255(2cm)=(180,215,238)
                rgb255(3cm)=(135,188,222)
                rgb255(4cm)=(90,158,205)
                rgb255(6cm)=(50,125,182)
                rgb255(8cm)=(25,95,160)
                rgb255(10cm)=(15,70,140)
              },
            point meta min=0,
            point meta max=0.9,
            major tick length=0pt,
            unbounded coords=jump,
            axis background/.style={pattern=sparse north east lines, pattern color=black!25},
            enlargelimits=false,
            scale only axis,
            axis on top,
            ymin=0.5, ymax=16.5,
            ytick={1,2,...,16},
            y dir=reverse,
          }
      }

      \begin{axis}[
          heatbase,
          name=fa,
          title={Full AttnRes, Pre-Attn},
          ylabel={Layer},
          xlabel={Source Index},
          xmin=-0.5, xmax=30.5,
          xtick={0,5,10,15,20,25,30},
          width=320pt, height=160pt,
        ]
        \addplot[matrix plot*, mesh/cols=31, point meta=explicit] coordinates {
            (0,1) [1.000] (1,1) [nan] (2,1) [nan] (3,1) [nan] (4,1) [nan] (5,1) [nan] (6,1) [nan] (7,1) [nan] (8,1) [nan] (9,1) [nan] (10,1) [nan] (11,1) [nan] (12,1) [nan] (13,1) [nan] (14,1) [nan] (15,1) [nan] (16,1) [nan] (17,1) [nan] (18,1) [nan] (19,1) [nan] (20,1) [nan] (21,1) [nan] (22,1) [nan] (23,1) [nan] (24,1) [nan] (25,1) [nan] (26,1) [nan] (27,1) [nan] (28,1) [nan] (29,1) [nan] (30,1) [nan]
            (0,2) [0.291] (1,2) [0.320] (2,2) [0.389] (3,2) [nan] (4,2) [nan] (5,2) [nan] (6,2) [nan] (7,2) [nan] (8,2) [nan] (9,2) [nan] (10,2) [nan] (11,2) [nan] (12,2) [nan] (13,2) [nan] (14,2) [nan] (15,2) [nan] (16,2) [nan] (17,2) [nan] (18,2) [nan] (19,2) [nan] (20,2) [nan] (21,2) [nan] (22,2) [nan] (23,2) [nan] (24,2) [nan] (25,2) [nan] (26,2) [nan] (27,2) [nan] (28,2) [nan] (29,2) [nan] (30,2) [nan]
            (0,3) [0.145] (1,3) [0.132] (2,3) [0.109] (3,3) [0.348] (4,3) [0.266] (5,3) [nan] (6,3) [nan] (7,3) [nan] (8,3) [nan] (9,3) [nan] (10,3) [nan] (11,3) [nan] (12,3) [nan] (13,3) [nan] (14,3) [nan] (15,3) [nan] (16,3) [nan] (17,3) [nan] (18,3) [nan] (19,3) [nan] (20,3) [nan] (21,3) [nan] (22,3) [nan] (23,3) [nan] (24,3) [nan] (25,3) [nan] (26,3) [nan] (27,3) [nan] (28,3) [nan] (29,3) [nan] (30,3) [nan]
            (0,4) [0.041] (1,4) [0.079] (2,4) [0.025] (3,4) [0.030] (4,4) [0.034] (5,4) [0.755] (6,4) [0.037] (7,4) [nan] (8,4) [nan] (9,4) [nan] (10,4) [nan] (11,4) [nan] (12,4) [nan] (13,4) [nan] (14,4) [nan] (15,4) [nan] (16,4) [nan] (17,4) [nan] (18,4) [nan] (19,4) [nan] (20,4) [nan] (21,4) [nan] (22,4) [nan] (23,4) [nan] (24,4) [nan] (25,4) [nan] (26,4) [nan] (27,4) [nan] (28,4) [nan] (29,4) [nan] (30,4) [nan]
            (0,5) [0.087] (1,5) [0.025] (2,5) [0.042] (3,5) [0.079] (4,5) [0.091] (5,5) [0.141] (6,5) [0.131] (7,5) [0.120] (8,5) [0.283] (9,5) [nan] (10,5) [nan] (11,5) [nan] (12,5) [nan] (13,5) [nan] (14,5) [nan] (15,5) [nan] (16,5) [nan] (17,5) [nan] (18,5) [nan] (19,5) [nan] (20,5) [nan] (21,5) [nan] (22,5) [nan] (23,5) [nan] (24,5) [nan] (25,5) [nan] (26,5) [nan] (27,5) [nan] (28,5) [nan] (29,5) [nan] (30,5) [nan]
            (0,6) [0.049] (1,6) [0.010] (2,6) [0.013] (3,6) [0.011] (4,6) [0.016] (5,6) [0.022] (6,6) [0.015] (7,6) [0.032] (8,6) [0.044] (9,6) [0.516] (10,6) [0.274] (11,6) [nan] (12,6) [nan] (13,6) [nan] (14,6) [nan] (15,6) [nan] (16,6) [nan] (17,6) [nan] (18,6) [nan] (19,6) [nan] (20,6) [nan] (21,6) [nan] (22,6) [nan] (23,6) [nan] (24,6) [nan] (25,6) [nan] (26,6) [nan] (27,6) [nan] (28,6) [nan] (29,6) [nan] (30,6) [nan]
            (0,7) [0.055] (1,7) [0.012] (2,7) [0.014] (3,7) [0.013] (4,7) [0.017] (5,7) [0.019] (6,7) [0.016] (7,7) [0.019] (8,7) [0.050] (9,7) [0.083] (10,7) [0.084] (11,7) [0.371] (12,7) [0.247] (13,7) [nan] (14,7) [nan] (15,7) [nan] (16,7) [nan] (17,7) [nan] (18,7) [nan] (19,7) [nan] (20,7) [nan] (21,7) [nan] (22,7) [nan] (23,7) [nan] (24,7) [nan] (25,7) [nan] (26,7) [nan] (27,7) [nan] (28,7) [nan] (29,7) [nan] (30,7) [nan]
            (0,8) [0.022] (1,8) [0.014] (2,8) [0.009] (3,8) [0.007] (4,8) [0.007] (5,8) [0.189] (6,8) [0.004] (7,8) [0.021] (8,8) [0.012] (9,8) [0.132] (10,8) [0.027] (11,8) [0.133] (12,8) [0.029] (13,8) [0.337] (14,8) [0.057] (15,8) [nan] (16,8) [nan] (17,8) [nan] (18,8) [nan] (19,8) [nan] (20,8) [nan] (21,8) [nan] (22,8) [nan] (23,8) [nan] (24,8) [nan] (25,8) [nan] (26,8) [nan] (27,8) [nan] (28,8) [nan] (29,8) [nan] (30,8) [nan]
            (0,9) [0.058] (1,9) [0.055] (2,9) [0.036] (3,9) [0.019] (4,9) [0.023] (5,9) [0.015] (6,9) [0.015] (7,9) [0.006] (8,9) [0.035] (9,9) [0.047] (10,9) [0.052] (11,9) [0.091] (12,9) [0.072] (13,9) [0.078] (14,9) [0.154] (15,9) [0.065] (16,9) [0.181] (17,9) [nan] (18,9) [nan] (19,9) [nan] (20,9) [nan] (21,9) [nan] (22,9) [nan] (23,9) [nan] (24,9) [nan] (25,9) [nan] (26,9) [nan] (27,9) [nan] (28,9) [nan] (29,9) [nan] (30,9) [nan]
            (0,10) [0.049] (1,10) [0.042] (2,10) [0.054] (3,10) [0.028] (4,10) [0.022] (5,10) [0.011] (6,10) [0.017] (7,10) [0.004] (8,10) [0.031] (9,10) [0.045] (10,10) [0.071] (11,10) [0.126] (12,10) [0.060] (13,10) [0.024] (14,10) [0.059] (15,10) [0.031] (16,10) [0.050] (17,10) [0.169] (18,10) [0.107] (19,10) [nan] (20,10) [nan] (21,10) [nan] (22,10) [nan] (23,10) [nan] (24,10) [nan] (25,10) [nan] (26,10) [nan] (27,10) [nan] (28,10) [nan] (29,10) [nan] (30,10) [nan]
            (0,11) [0.036] (1,11) [0.056] (2,11) [0.136] (3,11) [0.019] (4,11) [0.022] (5,11) [0.033] (6,11) [0.005] (7,11) [0.007] (8,11) [0.008] (9,11) [0.019] (10,11) [0.057] (11,11) [0.109] (12,11) [0.038] (13,11) [0.016] (14,11) [0.046] (15,11) [0.028] (16,11) [0.038] (17,11) [0.126] (18,11) [0.089] (19,11) [0.066] (20,11) [0.047] (21,11) [nan] (22,11) [nan] (23,11) [nan] (24,11) [nan] (25,11) [nan] (26,11) [nan] (27,11) [nan] (28,11) [nan] (29,11) [nan] (30,11) [nan]
            (0,12) [0.021] (1,12) [0.015] (2,12) [0.014] (3,12) [0.007] (4,12) [0.013] (5,12) [0.152] (6,12) [0.004] (7,12) [0.017] (8,12) [0.021] (9,12) [0.024] (10,12) [0.021] (11,12) [0.052] (12,12) [0.018] (13,12) [0.092] (14,12) [0.039] (15,12) [0.028] (16,12) [0.032] (17,12) [0.254] (18,12) [0.047] (19,12) [0.062] (20,12) [0.029] (21,12) [0.014] (22,12) [0.024] (23,12) [nan] (24,12) [nan] (25,12) [nan] (26,12) [nan] (27,12) [nan] (28,12) [nan] (29,12) [nan] (30,12) [nan]
            (0,13) [0.051] (1,13) [0.023] (2,13) [0.033] (3,13) [0.031] (4,13) [0.025] (5,13) [0.012] (6,13) [0.014] (7,13) [0.004] (8,13) [0.040] (9,13) [0.017] (10,13) [0.017] (11,13) [0.040] (12,13) [0.031] (13,13) [0.025] (14,13) [0.042] (15,13) [0.008] (16,13) [0.064] (17,13) [0.101] (18,13) [0.073] (19,13) [0.070] (20,13) [0.063] (21,13) [0.005] (22,13) [0.075] (23,13) [0.052] (24,13) [0.082] (25,13) [nan] (26,13) [nan] (27,13) [nan] (28,13) [nan] (29,13) [nan] (30,13) [nan]
            (0,14) [0.036] (1,14) [0.017] (2,14) [0.015] (3,14) [0.010] (4,14) [0.023] (5,14) [0.092] (6,14) [0.003] (7,14) [0.047] (8,14) [0.018] (9,14) [0.017] (10,14) [0.038] (11,14) [0.083] (12,14) [0.033] (13,14) [0.021] (14,14) [0.039] (15,14) [0.045] (16,14) [0.029] (17,14) [0.153] (18,14) [0.044] (19,14) [0.045] (20,14) [0.026] (21,14) [0.005] (22,14) [0.035] (23,14) [0.062] (24,14) [0.024] (25,14) [0.012] (26,14) [0.031] (27,14) [nan] (28,14) [nan] (29,14) [nan] (30,14) [nan]
            (0,15) [0.032] (1,15) [0.030] (2,15) [0.095] (3,15) [0.050] (4,15) [0.024] (5,15) [0.011] (6,15) [0.005] (7,15) [0.007] (8,15) [0.017] (9,15) [0.006] (10,15) [0.020] (11,15) [0.048] (12,15) [0.019] (13,15) [0.006] (14,15) [0.013] (15,15) [0.012] (16,15) [0.018] (17,15) [0.014] (18,15) [0.024] (19,15) [0.051] (20,15) [0.013] (21,15) [0.003] (22,15) [0.015] (23,15) [0.012] (24,15) [0.008] (25,15) [0.011] (26,15) [0.007] (27,15) [0.417] (28,15) [0.011] (29,15) [nan] (30,15) [nan]
            (0,16) [0.012] (1,16) [0.145] (2,16) [0.043] (3,16) [0.006] (4,16) [0.006] (5,16) [0.087] (6,16) [0.002] (7,16) [0.002] (8,16) [0.002] (9,16) [0.005] (10,16) [0.007] (11,16) [0.012] (12,16) [0.003] (13,16) [0.012] (14,16) [0.008] (15,16) [0.002] (16,16) [0.004] (17,16) [0.016] (18,16) [0.009] (19,16) [0.001] (20,16) [0.003] (21,16) [0.002] (22,16) [0.003] (23,16) [0.005] (24,16) [0.001] (25,16) [0.001] (26,16) [0.001] (27,16) [0.596] (28,16) [0.001] (29,16) [0.002] (30,16) [0.001]
          };
      \end{axis}

      \begin{axis}[
          heatbase,
          name=fm,
          at={($(fa.north east)+(0.5cm,0)$)},
          anchor=north west,
          title={Full AttnRes, Pre-MLP},
          ylabel={},
          yticklabels={},
          xlabel={Source Index},
          xmin=-0.5, xmax=31.5,
          xtick={0,5,10,15,20,25,30},
          width=320pt, height=160pt,
        ]
        \addplot[matrix plot*, mesh/cols=32, point meta=explicit] coordinates {
            (0,1) [0.507] (1,1) [0.493] (2,1) [nan] (3,1) [nan] (4,1) [nan] (5,1) [nan] (6,1) [nan] (7,1) [nan] (8,1) [nan] (9,1) [nan] (10,1) [nan] (11,1) [nan] (12,1) [nan] (13,1) [nan] (14,1) [nan] (15,1) [nan] (16,1) [nan] (17,1) [nan] (18,1) [nan] (19,1) [nan] (20,1) [nan] (21,1) [nan] (22,1) [nan] (23,1) [nan] (24,1) [nan] (25,1) [nan] (26,1) [nan] (27,1) [nan] (28,1) [nan] (29,1) [nan] (30,1) [nan] (31,1) [nan]
            (0,2) [0.107] (1,2) [0.138] (2,2) [0.051] (3,2) [0.704] (4,2) [nan] (5,2) [nan] (6,2) [nan] (7,2) [nan] (8,2) [nan] (9,2) [nan] (10,2) [nan] (11,2) [nan] (12,2) [nan] (13,2) [nan] (14,2) [nan] (15,2) [nan] (16,2) [nan] (17,2) [nan] (18,2) [nan] (19,2) [nan] (20,2) [nan] (21,2) [nan] (22,2) [nan] (23,2) [nan] (24,2) [nan] (25,2) [nan] (26,2) [nan] (27,2) [nan] (28,2) [nan] (29,2) [nan] (30,2) [nan] (31,2) [nan]
            (0,3) [0.056] (1,3) [0.127] (2,3) [0.035] (3,3) [0.406] (4,3) [0.036] (5,3) [0.340] (6,3) [nan] (7,3) [nan] (8,3) [nan] (9,3) [nan] (10,3) [nan] (11,3) [nan] (12,3) [nan] (13,3) [nan] (14,3) [nan] (15,3) [nan] (16,3) [nan] (17,3) [nan] (18,3) [nan] (19,3) [nan] (20,3) [nan] (21,3) [nan] (22,3) [nan] (23,3) [nan] (24,3) [nan] (25,3) [nan] (26,3) [nan] (27,3) [nan] (28,3) [nan] (29,3) [nan] (30,3) [nan] (31,3) [nan]
            (0,4) [0.058] (1,4) [0.101] (2,4) [0.035] (3,4) [0.373] (4,4) [0.035] (5,4) [0.263] (6,4) [0.080] (7,4) [0.054] (8,4) [nan] (9,4) [nan] (10,4) [nan] (11,4) [nan] (12,4) [nan] (13,4) [nan] (14,4) [nan] (15,4) [nan] (16,4) [nan] (17,4) [nan] (18,4) [nan] (19,4) [nan] (20,4) [nan] (21,4) [nan] (22,4) [nan] (23,4) [nan] (24,4) [nan] (25,4) [nan] (26,4) [nan] (27,4) [nan] (28,4) [nan] (29,4) [nan] (30,4) [nan] (31,4) [nan]
            (0,5) [0.031] (1,5) [0.027] (2,5) [0.010] (3,5) [0.042] (4,5) [0.011] (5,5) [0.022] (6,5) [0.018] (7,5) [0.017] (8,5) [0.025] (9,5) [0.796] (10,5) [nan] (11,5) [nan] (12,5) [nan] (13,5) [nan] (14,5) [nan] (15,5) [nan] (16,5) [nan] (17,5) [nan] (18,5) [nan] (19,5) [nan] (20,5) [nan] (21,5) [nan] (22,5) [nan] (23,5) [nan] (24,5) [nan] (25,5) [nan] (26,5) [nan] (27,5) [nan] (28,5) [nan] (29,5) [nan] (30,5) [nan] (31,5) [nan]
            (0,6) [0.032] (1,6) [0.042] (2,6) [0.010] (3,6) [0.029] (4,6) [0.006] (5,6) [0.014] (6,6) [0.012] (7,6) [0.007] (8,6) [0.016] (9,6) [0.242] (10,6) [0.041] (11,6) [0.549] (12,6) [nan] (13,6) [nan] (14,6) [nan] (15,6) [nan] (16,6) [nan] (17,6) [nan] (18,6) [nan] (19,6) [nan] (20,6) [nan] (21,6) [nan] (22,6) [nan] (23,6) [nan] (24,6) [nan] (25,6) [nan] (26,6) [nan] (27,6) [nan] (28,6) [nan] (29,6) [nan] (30,6) [nan] (31,6) [nan]
            (0,7) [0.028] (1,7) [0.043] (2,7) [0.010] (3,7) [0.016] (4,7) [0.009] (5,7) [0.009] (6,7) [0.012] (7,7) [0.007] (8,7) [0.017] (9,7) [0.112] (10,7) [0.025] (11,7) [0.126] (12,7) [0.046] (13,7) [0.539] (14,7) [nan] (15,7) [nan] (16,7) [nan] (17,7) [nan] (18,7) [nan] (19,7) [nan] (20,7) [nan] (21,7) [nan] (22,7) [nan] (23,7) [nan] (24,7) [nan] (25,7) [nan] (26,7) [nan] (27,7) [nan] (28,7) [nan] (29,7) [nan] (30,7) [nan] (31,7) [nan]
            (0,8) [0.029] (1,8) [0.106] (2,8) [0.026] (3,8) [0.065] (4,8) [0.011] (5,8) [0.028] (6,8) [0.016] (7,8) [0.002] (8,8) [0.016] (9,8) [0.100] (10,8) [0.010] (11,8) [0.070] (12,8) [0.021] (13,8) [0.373] (14,8) [0.084] (15,8) [0.043] (16,8) [nan] (17,8) [nan] (18,8) [nan] (19,8) [nan] (20,8) [nan] (21,8) [nan] (22,8) [nan] (23,8) [nan] (24,8) [nan] (25,8) [nan] (26,8) [nan] (27,8) [nan] (28,8) [nan] (29,8) [nan] (30,8) [nan] (31,8) [nan]
            (0,9) [0.028] (1,9) [0.121] (2,9) [0.021] (3,9) [0.048] (4,9) [0.010] (5,9) [0.013] (6,9) [0.012] (7,9) [0.003] (8,9) [0.012] (9,9) [0.062] (10,9) [0.010] (11,9) [0.053] (12,9) [0.017] (13,9) [0.099] (14,9) [0.042] (15,9) [0.024] (16,9) [0.041] (17,9) [0.384] (18,9) [nan] (19,9) [nan] (20,9) [nan] (21,9) [nan] (22,9) [nan] (23,9) [nan] (24,9) [nan] (25,9) [nan] (26,9) [nan] (27,9) [nan] (28,9) [nan] (29,9) [nan] (30,9) [nan] (31,9) [nan]
            (0,10) [0.021] (1,10) [0.152] (2,10) [0.020] (3,10) [0.047] (4,10) [0.006] (5,10) [0.024] (6,10) [0.011] (7,10) [0.002] (8,10) [0.008] (9,10) [0.057] (10,10) [0.005] (11,10) [0.030] (12,10) [0.008] (13,10) [0.095] (14,10) [0.026] (15,10) [0.011] (16,10) [0.026] (17,10) [0.249] (18,10) [0.048] (19,10) [0.156] (20,10) [nan] (21,10) [nan] (22,10) [nan] (23,10) [nan] (24,10) [nan] (25,10) [nan] (26,10) [nan] (27,10) [nan] (28,10) [nan] (29,10) [nan] (30,10) [nan] (31,10) [nan]
            (0,11) [0.021] (1,11) [0.095] (2,11) [0.012] (3,11) [0.031] (4,11) [0.005] (5,11) [0.011] (6,11) [0.009] (7,11) [0.005] (8,11) [0.008] (9,11) [0.106] (10,11) [0.011] (11,11) [0.070] (12,11) [0.014] (13,11) [0.079] (14,11) [0.025] (15,11) [0.020] (16,11) [0.014] (17,11) [0.191] (18,11) [0.029] (19,11) [0.217] (20,11) [0.024] (21,11) [0.003] (22,11) [nan] (23,11) [nan] (24,11) [nan] (25,11) [nan] (26,11) [nan] (27,11) [nan] (28,11) [nan] (29,11) [nan] (30,11) [nan] (31,11) [nan]
            (0,12) [0.023] (1,12) [0.147] (2,12) [0.025] (3,12) [0.061] (4,12) [0.007] (5,12) [0.022] (6,12) [0.012] (7,12) [0.015] (8,12) [0.009] (9,12) [0.064] (10,12) [0.005] (11,12) [0.033] (12,12) [0.007] (13,12) [0.091] (14,12) [0.020] (15,12) [0.011] (16,12) [0.021] (17,12) [0.167] (18,12) [0.029] (19,12) [0.123] (20,12) [0.026] (21,12) [0.009] (22,12) [0.035] (23,12) [0.040] (24,12) [nan] (25,12) [nan] (26,12) [nan] (27,12) [nan] (28,12) [nan] (29,12) [nan] (30,12) [nan] (31,12) [nan]
            (0,13) [0.021] (1,13) [0.135] (2,13) [0.031] (3,13) [0.025] (4,13) [0.008] (5,13) [0.015] (6,13) [0.010] (7,13) [0.003] (8,13) [0.009] (9,13) [0.032] (10,13) [0.004] (11,13) [0.018] (12,13) [0.007] (13,13) [0.044] (14,13) [0.017] (15,13) [0.005] (16,13) [0.022] (17,13) [0.109] (18,13) [0.029] (19,13) [0.077] (20,13) [0.029] (21,13) [0.006] (22,13) [0.033] (23,13) [0.027] (24,13) [0.033] (25,13) [0.249] (26,13) [nan] (27,13) [nan] (28,13) [nan] (29,13) [nan] (30,13) [nan] (31,13) [nan]
            (0,14) [0.020] (1,14) [0.107] (2,14) [0.021] (3,14) [0.036] (4,14) [0.007] (5,14) [0.008] (6,14) [0.008] (7,14) [0.003] (8,14) [0.007] (9,14) [0.038] (10,14) [0.005] (11,14) [0.027] (12,14) [0.008] (13,14) [0.033] (14,14) [0.014] (15,14) [0.008] (16,14) [0.013] (17,14) [0.073] (18,14) [0.019] (19,14) [0.151] (20,14) [0.018] (21,14) [0.007] (22,14) [0.038] (23,14) [0.019] (24,14) [0.023] (25,14) [0.179] (26,14) [0.033] (27,14) [0.078] (28,14) [nan] (29,14) [nan] (30,14) [nan] (31,14) [nan]
            (0,15) [0.025] (1,15) [0.063] (2,15) [0.017] (3,15) [0.031] (4,15) [0.006] (5,15) [0.007] (6,15) [0.006] (7,15) [0.003] (8,15) [0.005] (9,15) [0.017] (10,15) [0.003] (11,15) [0.015] (12,15) [0.005] (13,15) [0.023] (14,15) [0.010] (15,15) [0.006] (16,15) [0.015] (17,15) [0.039] (18,15) [0.020] (19,15) [0.115] (20,15) [0.022] (21,15) [0.039] (22,15) [0.034] (23,15) [0.024] (24,15) [0.028] (25,15) [0.185] (26,15) [0.047] (27,15) [0.061] (28,15) [0.049] (29,15) [0.078] (30,15) [nan] (31,15) [nan]
            (0,16) [0.026] (1,16) [0.021] (2,16) [0.010] (3,16) [0.017] (4,16) [0.003] (5,16) [0.006] (6,16) [0.005] (7,16) [0.002] (8,16) [0.004] (9,16) [0.013] (10,16) [0.002] (11,16) [0.012] (12,16) [0.004] (13,16) [0.017] (14,16) [0.005] (15,16) [0.003] (16,16) [0.009] (17,16) [0.015] (18,16) [0.015] (19,16) [0.068] (20,16) [0.019] (21,16) [0.072] (22,16) [0.033] (23,16) [0.009] (24,16) [0.026] (25,16) [0.089] (26,16) [0.041] (27,16) [0.037] (28,16) [0.047] (29,16) [0.121] (30,16) [0.063] (31,16) [0.185]
          };
      \end{axis}

      \begin{axis}[
          heatbase,
          name=ba,
          at={($(fa.south west)-(0,2.0cm)$)},
          anchor=north west,
          title={Block AttnRes, Pre-Attn},
          ylabel={Layer},
          xlabel={Block Index},
          xmin=-0.5, xmax=8.5,
          xtick={0,1,...,8},
          width=320pt, height=160pt,
        ]
        \addplot[matrix plot*, mesh/cols=9, point meta=explicit] coordinates {
            (0,1) [1.000] (1,1) [nan] (2,1) [nan] (3,1) [nan] (4,1) [nan] (5,1) [nan] (6,1) [nan] (7,1) [nan] (8,1) [nan]
            (0,2) [0.764] (1,2) [0.236] (2,2) [nan] (3,2) [nan] (4,2) [nan] (5,2) [nan] (6,2) [nan] (7,2) [nan] (8,2) [nan]
            (0,3) [0.729] (1,3) [0.271] (2,3) [nan] (3,3) [nan] (4,3) [nan] (5,3) [nan] (6,3) [nan] (7,3) [nan] (8,3) [nan]
            (0,4) [0.210] (1,4) [0.417] (2,4) [0.373] (3,4) [nan] (4,4) [nan] (5,4) [nan] (6,4) [nan] (7,4) [nan] (8,4) [nan]
            (0,5) [0.228] (1,5) [0.070] (2,5) [0.702] (3,5) [nan] (4,5) [nan] (5,5) [nan] (6,5) [nan] (7,5) [nan] (8,5) [nan]
            (0,6) [0.118] (1,6) [0.026] (2,6) [0.060] (3,6) [0.796] (4,6) [nan] (5,6) [nan] (6,6) [nan] (7,6) [nan] (8,6) [nan]
            (0,7) [0.204] (1,7) [0.027] (2,7) [0.122] (3,7) [0.647] (4,7) [nan] (5,7) [nan] (6,7) [nan] (7,7) [nan] (8,7) [nan]
            (0,8) [0.130] (1,8) [0.061] (2,8) [0.052] (3,8) [0.382] (4,8) [0.376] (5,8) [nan] (6,8) [nan] (7,8) [nan] (8,8) [nan]
            (0,9) [0.151] (1,9) [0.018] (2,9) [0.079] (3,9) [0.222] (4,9) [0.530] (5,9) [nan] (6,9) [nan] (7,9) [nan] (8,9) [nan]
            (0,10) [0.089] (1,10) [0.011] (2,10) [0.041] (3,10) [0.077] (4,10) [0.108] (5,10) [0.674] (6,10) [nan] (7,10) [nan] (8,10) [nan]
            (0,11) [0.152] (1,11) [0.116] (2,11) [0.071] (3,11) [0.122] (4,11) [0.132] (5,11) [0.407] (6,11) [nan] (7,11) [nan] (8,11) [nan]
            (0,12) [0.110] (1,12) [0.174] (2,12) [0.028] (3,12) [0.204] (4,12) [0.059] (5,12) [0.316] (6,12) [0.109] (7,12) [nan] (8,12) [nan]
            (0,13) [0.093] (1,13) [0.046] (2,13) [0.158] (3,13) [0.137] (4,13) [0.114] (5,13) [0.388] (6,13) [0.066] (7,13) [nan] (8,13) [nan]
            (0,14) [0.098] (1,14) [0.147] (2,14) [0.087] (3,14) [0.082] (4,14) [0.079] (5,14) [0.258] (6,14) [0.040] (7,14) [0.209] (8,14) [nan]
            (0,15) [0.088] (1,15) [0.075] (2,15) [0.056] (3,15) [0.054] (4,15) [0.051] (5,15) [0.158] (6,15) [0.299] (7,15) [0.219] (8,15) [nan]
            (0,16) [0.113] (1,16) [0.217] (2,16) [0.021] (3,16) [0.213] (4,16) [0.028] (5,16) [0.289] (6,16) [0.058] (7,16) [0.029] (8,16) [0.032]
          };
      \end{axis}

      \begin{axis}[
          heatbase,
          name=bm,
          at={($(ba.north east)+(0.5cm,0)$)},
          anchor=north west,
          title={Block AttnRes, Pre-MLP},
          ylabel={},
          yticklabels={},
          xlabel={Block Index},
          xmin=-0.5, xmax=8.5,
          xtick={0,1,...,8},
          width=320pt, height=160pt,
        ]
        \addplot[matrix plot*, mesh/cols=9, point meta=explicit] coordinates {
            (0,1) [0.510] (1,1) [0.490] (2,1) [nan] (3,1) [nan] (4,1) [nan] (5,1) [nan] (6,1) [nan] (7,1) [nan] (8,1) [nan]
            (0,2) [0.912] (1,2) [0.088] (2,2) [nan] (3,2) [nan] (4,2) [nan] (5,2) [nan] (6,2) [nan] (7,2) [nan] (8,2) [nan]
            (0,3) [0.448] (1,3) [0.086] (2,3) [0.465] (3,3) [nan] (4,3) [nan] (5,3) [nan] (6,3) [nan] (7,3) [nan] (8,3) [nan]
            (0,4) [0.557] (1,4) [0.147] (2,4) [0.296] (3,4) [nan] (4,4) [nan] (5,4) [nan] (6,4) [nan] (7,4) [nan] (8,4) [nan]
            (0,5) [0.080] (1,5) [0.020] (2,5) [0.035] (3,5) [0.866] (4,5) [nan] (5,5) [nan] (6,5) [nan] (7,5) [nan] (8,5) [nan]
            (0,6) [0.132] (1,6) [0.062] (2,6) [0.088] (3,6) [0.718] (4,6) [nan] (5,6) [nan] (6,6) [nan] (7,6) [nan] (8,6) [nan]
            (0,7) [0.100] (1,7) [0.029] (2,7) [0.038] (3,7) [0.102] (4,7) [0.731] (5,7) [nan] (6,7) [nan] (7,7) [nan] (8,7) [nan]
            (0,8) [0.107] (1,8) [0.077] (2,8) [0.061] (3,8) [0.147] (4,8) [0.607] (5,8) [nan] (6,8) [nan] (7,8) [nan] (8,8) [nan]
            (0,9) [0.056] (1,9) [0.013] (2,9) [0.013] (3,9) [0.024] (4,9) [0.061] (5,9) [0.832] (6,9) [nan] (7,9) [nan] (8,9) [nan]
            (0,10) [0.086] (1,10) [0.052] (2,10) [0.043] (3,10) [0.036] (4,10) [0.147] (5,10) [0.637] (6,10) [nan] (7,10) [nan] (8,10) [nan]
            (0,11) [0.075] (1,11) [0.053] (2,11) [0.023] (3,11) [0.019] (4,11) [0.100] (5,11) [0.254] (6,11) [0.476] (7,11) [nan] (8,11) [nan]
            (0,12) [0.060] (1,12) [0.134] (2,12) [0.033] (3,12) [0.028] (4,12) [0.093] (5,12) [0.224] (6,12) [0.427] (7,12) [nan] (8,12) [nan]
            (0,13) [0.059] (1,13) [0.130] (2,13) [0.021] (3,13) [0.026] (4,13) [0.064] (5,13) [0.131] (6,13) [0.129] (7,13) [0.439] (8,13) [nan]
            (0,14) [0.105] (1,14) [0.081] (2,14) [0.015] (3,14) [0.017] (4,14) [0.056] (5,14) [0.152] (6,14) [0.281] (7,14) [0.293] (8,14) [nan]
            (0,15) [0.085] (1,15) [0.057] (2,15) [0.014] (3,15) [0.012] (4,15) [0.033] (5,15) [0.088] (6,15) [0.268] (7,15) [0.233] (8,15) [0.211]
            (0,16) [0.093] (1,16) [0.068] (2,16) [0.018] (3,16) [0.014] (4,16) [0.033] (5,16) [0.082] (6,16) [0.240] (7,16) [0.234] (8,16) [0.218]
          };
      \end{axis}

      \begin{axis}[
          hide axis,
          scale only axis,
          width=0pt,
          height=340pt,
          at={($(fm.north east)!0.5!(bm.south east)+(1.0cm,0)$)},
          anchor=center,
          colormap={whiteblue}{
              rgb255(0cm)=(255,255,255)
              rgb255(1cm)=(220,235,248)
              rgb255(2cm)=(180,215,238)
              rgb255(3cm)=(135,188,222)
              rgb255(4cm)=(90,158,205)
              rgb255(6cm)=(50,125,182)
              rgb255(8cm)=(25,95,160)
              rgb255(10cm)=(15,70,140)
            },
          point meta min=0,
          point meta max=0.9,
          colorbar right,
          colorbar style={
              width=6pt,
              ytick={0, 0.2, 0.4, 0.6, 0.8},
              yticklabel style={font=\scriptsize},
              title={Weight},
              title style={font=\scriptsize, at={(0.5,1.02)}, anchor=south},
            },
        ]
        \addplot[draw=none, point meta=explicit] coordinates {(0,0) [0]};
      \end{axis}

    \end{tikzpicture}%
  }% end adjustbox
  \caption{
    Depth-wise attention weight distributions for a 16-head model with full (top) and block (bottom) Attention Residuals, averaged over tokens. The model has 16 attention and 16 MLP layers. Each row shows how the $l$th attention (left) or MLP (right) layer distributes weight over previous sources. Diagonal dominance indicates locality remains the primary information pathway, while persistent weights on source 0 (embedding) and occasional off-diagonal concentrations reveal learned skip connections. Block attention ($N=8$) recovers the essential structure with sharper, more decisive weight distributions.
  }
  \label{fig:attn-res-weights}
\end{figure*}


\label{sec:analysis}
 We visualize the learned weights $\brickred{\alpha_{i \to l}}$ in Fig.~\ref{fig:attn-res-weights} for the 16-head model (from Table~\ref{tab:baseline_mhc_attnres_compact}) with both full and block ($N{=}8$) AttnRes.
Each heatmap shows how the $l$th attention or MLP layer (rows) allocates its attention over previous sources (columns), with pre-attention and pre-MLP layers shown separately.
We highlight three key observations:
\begin{itemize}[leftmargin=*]
    \item \textbf{Preserved locality.} Each layer attends most strongly to its immediate predecessor, yet selective off-diagonal concentrations emerge (e.g., layer 4 attending to early sources, layers 15--16 reaching back under the block setting), indicating learned skip connections beyond the standard residual path.
    \item \textbf{Layer specialization.} The embedding $\bm{h}_1$ retains non-trivial weight throughout, especially in pre-attention layers. Pre-MLP inputs show sharper diagonal reliance on recent representations, while pre-attention inputs maintain broader receptive fields, consistent with attention routing information across layers and MLPs operating locally.
    \item \textbf{Block AttnRes preserves structure.} Diagonal dominance, embedding persistence, and layer specialization all transfer from the full to the block variant, suggesting that block-wise compression acts as implicit regularization while preserving the essential information pathways.
\end{itemize}
